% =============================================================================
% chapter6.tex — Deployment and User Interface
% Chapter 6 of the Codara PFE Report
% =============================================================================

\chapter{Deployment and User Interface}
\label{ch:deployment}

% ---------------------------------------------------------------------------
\section{Introduction}
\label{sec:ch6-intro}
% ---------------------------------------------------------------------------

A system that works in a Jupyter notebook is not the same as a system that
works in production. The fifteen-week development sprint did not treat deployment
as an afterthought: the Docker Compose configuration, GitHub Actions pipeline,
and Azure Key Vault integration were built in Week 13 and evolved alongside the
pipeline logic. This chapter documents the gap between the core pipeline
described in Chapter~\ref{ch:implementation} and the production-ready service
that the host company's engineers can actually use.

The chapter is organised into three sections. Section~\ref{sec:backend-deployment}
covers the FastAPI backend structure, the three-service abstraction layer, the
Docker containerisation, the six-job GitHub Actions CI/CD pipeline, and the
Azure Key Vault integration for secrets management. Section~\ref{sec:frontend-deployment}
presents the React frontend: technology choices, key interface views, and the
end-to-end integration flow from file upload to download. The chapter closes
with a discussion of the deployment architecture's strengths and known
limitations.

% ---------------------------------------------------------------------------
\section{Backend Deployment}
\label{sec:backend-deployment}
% ---------------------------------------------------------------------------

\subsection{FastAPI Structure, Endpoints, and Background Tasks}
\label{subsec:fastapi-structure}

\textbf{The technology.} FastAPI~[1] is an asynchronous Python web framework
built on Starlette and Pydantic. Its key advantage for this project is
\emph{type-driven development}: all request and response bodies are declared as
Pydantic v2 models, from which FastAPI automatically generates OpenAPI
documentation, request validation, and JSON serialisation. Combined with its
native async support (built on \texttt{asyncio}), FastAPI allows the API server
to handle many concurrent polling requests without blocking while pipeline
coroutines execute in background tasks.

\textbf{How we used it.} The FastAPI application in
\texttt{backend/api/main.py} initialises in three phases:

\begin{enumerate}
  \item \textbf{Startup}: SQLAlchemy creates the SQLite database tables (if
    absent), seeds the default admin and user accounts (printing generated
    passwords to stdout), initialises the \texttt{PipelineQueueService} daemon
    consumer thread, and registers all six routers.
  \item \textbf{Runtime}: the ASGI server (uvicorn) accepts HTTP requests on
    port 8000. Routing is handled by six FastAPI routers: \texttt{auth},
    \texttt{conversions}, \texttt{knowledge\_base}, \texttt{admin},
    \texttt{analytics}, \texttt{notifications} + \texttt{settings}.
  \item \textbf{Shutdown}: the queue consumer thread is signalled to exit
    and the database session factory is closed.
\end{enumerate}

CORS is configured at startup to allow requests from
\texttt{localhost:8080} (Docker/nginx), \texttt{localhost:5173} (Vite dev server),
and \texttt{127.0.0.1:8080}. In production, the \texttt{FRONTEND\_URL}
environment variable adds the Azure Static Web App domain to the CORS
allowed-origins list.

[FIGURE 6.1 --- Caption: ``FastAPI router hierarchy: six routers mapped to
the main application, with authentication dependency injected at the router
level.'' --- Suggested: tree diagram with the FastAPI app at the root and six
router boxes as children, each annotated with its prefix and auth requirement.]

\subsection{Service Layer: Pipeline, Blob, and Queue}
\label{subsec:service-layer}

\textbf{The problem.} In the early sprints, the conversion router called the
LangGraph pipeline directly via \texttt{BackgroundTasks.add\_task()}. This
worked locally but created three problems at scale: (1) the pipeline
instantiated a new LangGraph graph object on every conversion request, adding
initialisation overhead; (2) file I/O was mixed with HTTP logic, making the
route handlers untestable without a running filesystem; (3) there was no way
to dispatch jobs to a remote worker process.

\textbf{The solution.} Three service modules were extracted in Week 15 to
decouple HTTP routing from infrastructure:

\textbf{\texttt{pipeline\_service.py} (\texttt{run\_pipeline\_sync})}: the
single function called by both the queue consumer and the BackgroundTasks
fallback. It downloads the SAS file from Blob Storage (or reads it from the
local \texttt{uploads/} directory), constructs the \texttt{PipelineState}
initial dict, executes the 8-node LangGraph pipeline synchronously (using
\texttt{asyncio.run()} inside the thread), updates the conversion record with
the result, and cleans up the temporary file. All pipeline errors are caught,
logged to DuckDB via \texttt{LLMAuditLogger}, and written to the
\texttt{conversion\_stages} table as failed stages.

\textbf{\texttt{blob\_service.py} (\texttt{BlobStorageService})}: provides
four methods: \texttt{upload\_file()} (write bytes to Azure Blob or local disk),
\texttt{download\_to\_temp()} (download blob to a named temp file, returning
the path), \texttt{list\_files()} (list blobs under a prefix), and
\texttt{delete\_folder()} (delete all blobs under a prefix). The service checks
for \texttt{AZURE\_STORAGE\_CONNECTION\_STRING} at initialisation; if absent,
all methods operate on the local \texttt{backend/uploads/} directory instead.
This transparent fallback means the same code path is exercised in both
development (local disk) and production (Azure Blob).

\textbf{\texttt{queue\_service.py} (\texttt{PipelineQueueService})}: wraps
Azure Queue Storage for asynchronous job dispatch. The \texttt{enqueue()} method
base64-encodes a JSON payload (containing the conversion ID and file paths) and
sends it to the Azure Queue with \texttt{visibility\_timeout=300} seconds. A
daemon thread started at application startup polls the queue every 2 seconds
via \texttt{receive\_messages(max\_messages=1, visibility\_timeout=300)} and
calls \texttt{run\_pipeline\_sync} for each message. The poison-message guard
increments a dequeue-count counter per message ID and dead-letters any message
with count $\geq 5$ to prevent infinite retry loops on unprocessable jobs.

[FIGURE 6.2 --- Caption: ``Service layer architecture: three services abstract
Azure Blob, Azure Queue, and the LangGraph pipeline from the HTTP routes.''
--- Suggested: three-column diagram with HTTP Routes on left, Services in
middle, and Azure/Pipeline on right, with arrows showing the call direction.]

\subsection{Containerisation and CI/CD}
\label{subsec:docker}

\textbf{The technology.} Docker~[2] enables reproducible, environment-independent
deployment by packaging the application and all its dependencies into an image.
Docker Compose~[3] orchestrates multiple containers (Redis, backend, frontend)
as a single service. GitHub Actions~[4] provides the CI/CD pipeline that
automates testing, building, and deploying on every push to main.

\textbf{Docker setup.} The backend image is built from a multi-stage
\texttt{infra/Dockerfile} based on \texttt{python:3.11-slim}. Stage 1 installs
build dependencies (gcc, libpq-dev) and builds the Python wheels; Stage 2
copies only the compiled wheels into a clean image, reducing the final image
size. The application listens on port 8000 and is started with uvicorn.

The \texttt{infra/docker-compose.yml} defines three services:
\texttt{redis} (image: \texttt{redis:7-alpine}, port 6379, no persistence),
\texttt{backend} (built from Dockerfile, port 8000, depends on redis),
\texttt{frontend} (built from the React app with nginx, port 8080, proxies
\texttt{/api} to the backend container).

\textbf{GitHub Actions pipeline.} The CI/CD pipeline comprises six jobs that
run on every push to the \texttt{main} branch:

[TABLE 6.1 --- Caption: ``GitHub Actions CI/CD pipeline: six jobs, their
triggers, and success criteria.'' --- Columns: Job, Trigger, Runtime, Success
Criterion.]

\begin{enumerate}
  \item \textbf{lint}: runs \texttt{flake8} and \texttt{mypy} on the backend
    codebase. Blocks merge on any type error or style violation above severity
    E/W.
  \item \textbf{test}: installs dev dependencies from
    \texttt{requirements/dev.txt} (which includes pytest via \texttt{base.txt}),
    runs \texttt{pytest tests/ -v --tb=short --cov=partition --cov-report=xml},
    and uploads the coverage report. Fails if any test fails.
  \item \textbf{benchmark}: runs the boundary detection benchmark
    (\texttt{benchmark/boundary\_benchmark.py}) and the oracle regression runner
    (\texttt{benchmark/regression\_runner.py}) against the gold-standard corpus.
    Fails if any gold-standard pair produces SEMANTIC\_FAIL.
  \item \textbf{docker-build}: builds both the backend and frontend Docker
    images without pushing. Catches Dockerfile errors and dependency conflicts.
  \item \textbf{codeql}: GitHub's default CodeQL security scanning for Python
    (CWE-089, CWE-022, CWE-078, and 12 other high-severity patterns). All
    CodeQL High findings were resolved in a dedicated security sprint (5
    alerts fixed across path traversal, SQL injection, and OS command injection
    categories).
  \item \textbf{azure-deploy}: on successful completion of all five preceding
    jobs, updates the Azure Container App image using:
    \texttt{az containerapp update --name codara-backend --resource-group rg-codara
    --image \$REGISTRY/backend:\$SHA}. The image-only update preserves all
    existing secrets and environment variable bindings, unlike
    \texttt{az containerapp up} which resets them on every deploy.
\end{enumerate}

\subsection{Azure Key Vault Integration}
\label{subsec:keyvault}

\textbf{The technology.} Azure Key Vault~[5] is a cloud secrets manager that
stores cryptographic keys, secrets, and certificates. Azure Container Apps can
reference Key Vault secrets directly as environment variables via
\texttt{--secrets keyvaultref:<URI>}, eliminating the need to store secrets
in GitHub Actions variables or \texttt{.env} files.

\textbf{How we used it.} The \texttt{infra/azure\_setup.sh} script provisions
six secrets in Key Vault:
\texttt{OLLAMA-API-KEY}, \texttt{OLLAMA-BASE-URL},
\texttt{GROQ-API-KEY}, \texttt{CODARA-JWT-SECRET},
\texttt{AZURE-OPENAI-API-KEY}, and \texttt{FRONTEND-URL} (for CORS).
The Container App's system-assigned managed identity is granted the
``Key Vault Secrets User'' role, allowing it to read secrets without any stored
credentials. The Container App environment also receives \texttt{CORS\_ORIGINS}
and \texttt{OLLAMA\_MODEL} as plain environment variables (non-sensitive).

This design satisfies NF-06 (JWT required on all endpoints), NF-07 (bcrypt
hashing), and NF-08 (rate limiting) at the infrastructure level: the JWT secret
is never stored in source code or container images, only fetched at runtime from
Key Vault.

% ---------------------------------------------------------------------------
\section{Frontend Dashboard}
\label{sec:frontend-deployment}
% ---------------------------------------------------------------------------

\subsection{Technology Stack and Architecture}
\label{subsec:frontend-stack}

\textbf{The technology.} The frontend is a React 18~[6] single-page application
built with Vite~[7] as the bundler, TypeScript~[8] for static typing, Tailwind
CSS~[9] + shadcn/ui~[10] for the component library, and Zustand~[11] for
global state management.

React 18 was chosen for its concurrent rendering capabilities (important for
the live stage-progress view) and its large ecosystem of UI libraries compatible
with Tailwind. Vite provides near-instant hot-module replacement in development
and produces optimised production bundles in under 10 seconds. TypeScript
enforces type safety across all API calls: the
\texttt{frontend/src/types/index.ts} file defines \texttt{Conversion},
\texttt{SasFile}, \texttt{PipelineStageInfo}, and related types that mirror
the Pydantic schemas on the backend.

Zustand was preferred over Redux for its minimal boilerplate: the
\texttt{conversion-store.ts} contains the full upload/start/poll/stop logic
in under 100 lines, with the polling interval hardcoded to 1,200 ms and
automatic stop when the conversion status is terminal.

All API calls pass through the \texttt{frontend/src/lib/api.ts} thin wrapper,
which reads the JWT token from \texttt{localStorage} under the key
\texttt{codara\_token} and attaches it as a Bearer header. The Vite dev-server
proxy forwards all \texttt{/api} requests to \texttt{http://localhost:8000},
so the frontend code never references the backend URL directly.

\subsection{Key Interface Views}
\label{subsec:workspace-ui}

The frontend exposes seven distinct views:

\textbf{Login / Register} (\texttt{src/pages/Login.tsx}): a minimal JWT
authentication form with error handling for invalid credentials and rate-limited
responses. On successful login, the JWT is stored in localStorage and the user
is redirected to the Dashboard.

\textbf{Dashboard} (\texttt{src/pages/Dashboard.tsx}): the primary landing
page. Lists the user's conversions with status badges (queued / running /
completed / partial / failed), accuracy scores where available, and creation
timestamps. Provides a ``New Conversion'' button that navigates to the Workspace.

[FIGURE 6.3 --- Caption: ``\codara{} Dashboard: list of conversions with
status badges, accuracy scores, and navigation to the Workspace.'' ---
Suggested: screenshot of the Dashboard page.]

\textbf{Workspace} (\texttt{src/pages/Workspace.tsx}): the main conversion
UI. Divided into three panels:
\begin{itemize}
  \item \textit{Upload panel}: drag-and-drop file upload with progress indicator.
    After upload, displays the file list with detected SAS structure metadata.
  \item \textit{Progress panel}: a live stage tracker showing the eight
    pipeline stages as a vertical list with status icons, latency figures, and
    warning counts. Stages update in real time as the 1,200 ms poll detects
    changes.
  \item \textit{Diff panel}: a GitHub-style side-by-side diff view
    (SAS source on the left, Python output on the right) using the
    \texttt{framer-motion} library for smooth transition animations. The SCS
    badge appears above each translated block, colour-coded by verdict
    (green = VERIFIED, blue = LIKELY\_CORRECT, yellow = UNCERTAIN,
    red = LIKELY\_INCORRECT).
\end{itemize}

[FIGURE 6.4 --- Caption: ``\codara{} Workspace: upload panel (left), live
stage tracker (centre), and GitHub-style SAS/Python side-by-side diff (right)
with SemantiCheck Score badges.'' --- Suggested: screenshot or mockup of the
three-panel Workspace view.]

\textbf{Admin views} (\texttt{src/pages/admin/}): six admin-only pages
accessible via the sidebar navigation for users with \texttt{role=admin}:
\textit{Users} (user management with role and status controls),
\textit{KBManagement} (knowledge-base CRUD with search),
\textit{KBChangelog} (version history of KB entries),
\textit{AuditLogs} (DuckDB-backed LLM cost and latency table),
\textit{SystemHealth} (Redis status, LanceDB table size, SQLite size,
circuit breaker states), and \textit{PipelineConfig} (editable pipeline
parameters including LLM timeouts and RAPTOR cluster count).

[FIGURE 6.5 --- Caption: ``Admin AuditLogs view: per-model latency and
estimated token cost, sortable by timestamp.'' --- Suggested: screenshot
of the AuditLogs admin page.]

\textbf{Knowledge Base} (\texttt{src/pages/KnowledgeBase.tsx}): semantic search
interface for the LanceDB KB. Users can search by SAS snippet or Python
translation text; results are displayed as SAS--Python pair cards with
confidence scores and complexity tier badges.

\subsection{End-to-End Integration Flow}
\label{subsec:e2e-integration}

The complete flow from the user's perspective is as follows:

\begin{enumerate}
  \item User navigates to the Workspace and drags a \texttt{.sas} file onto
    the upload zone. The file is sent via \texttt{POST /api/conversions/upload};
    the response includes a list of SAS file metadata objects.
  \item User clicks ``Start Conversion''. A \texttt{POST /api/conversions/start}
    request is sent with the file UUIDs and target runtime (\texttt{python}).
    The response returns a \texttt{conversion\_id}.
  \item The Zustand \texttt{conversion-store} begins polling
    \texttt{GET /api/conversions/\{id\}} every 1,200 ms. The stage tracker
    updates live as each of the eight pipeline stages completes.
  \item On completion, the diff panel renders the side-by-side view with SCS
    badges. The user can click any translated block to view the full
    SemantiCheck breakdown, Z3 verification outcome, and CDAIS test results.
  \item If any partition is LIKELY\_INCORRECT, the user can click ``Submit
    Correction'' to provide a corrected translation, which is ingested into
    the knowledge base via \texttt{POST /api/conversions/\{id\}/corrections}.
  \item The user clicks ``Download'' to receive a ZIP archive containing the
    merged Python script and the HTML conversion report.
\end{enumerate}

[FIGURE 6.6 --- Caption: ``End-to-end integration sequence diagram: browser,
FastAPI, Queue, Pipeline worker, and data stores.'' --- Suggested: UML
sequence diagram with five actors.]

% ---------------------------------------------------------------------------
\section{Conclusion}
\label{sec:ch6-conclusion}
% ---------------------------------------------------------------------------

This chapter demonstrated that \codara{} is not merely a research prototype
but a deployable production service. The FastAPI application handles
authentication, file upload, pipeline dispatch, and result polling through a
clean, layered API. The three service modules (pipeline, blob, queue) decouple
the HTTP layer from infrastructure, enabling both local development with
Docker Compose and cloud deployment on Azure Container Apps without code changes.
The six-job GitHub Actions pipeline provides continuous quality assurance:
every push to main triggers lint, unit tests, semantic regression tests,
Docker build validation, CodeQL security scanning, and an automated Azure deploy.

The React frontend provides a polished, production-ready user experience: the
three-panel Workspace view with live stage tracking and GitHub-style diff output
gives engineers the visibility they need to trust the system's translations,
while the SCS badges and one-click correction flow create a feedback loop that
continuously improves the knowledge base.

The quantitative performance of all these components against the targets
established in Chapter~\ref{ch:specifications} is measured in the following
chapter.

% ---------------------------------------------------------------------------
% References for Chapter 6
% ---------------------------------------------------------------------------
%
% [1]  Ramírez, S., "FastAPI," https://fastapi.tiangolo.com/, 2024.
% [2]  Docker Inc., "Docker: Empowering App Development for Developers,"
%      https://www.docker.com/, 2024.
% [3]  Docker Inc., "Docker Compose," https://docs.docker.com/compose/, 2024.
% [4]  GitHub Inc., "GitHub Actions,"
%      https://docs.github.com/en/actions, 2024.
% [5]  Microsoft, "Azure Key Vault,"
%      https://learn.microsoft.com/azure/key-vault/, 2024.
% [6]  Meta Open Source, "React 18," https://react.dev/, 2024.
% [7]  Evan You, "Vite: Next Generation Frontend Tooling,"
%      https://vitejs.dev/, 2024.
% [8]  Microsoft, "TypeScript," https://www.typescriptlang.org/, 2024.
% [9]  Tailwind Labs, "Tailwind CSS," https://tailwindcss.com/, 2024.
% [10] shadcn, "shadcn/ui," https://ui.shadcn.com/, 2024.
% [11] Pmndrs, "Zustand," https://github.com/pmndrs/zustand, 2024.
